% SEGMENT OLP-0096-S01
% Part: first-order-logic
% Chapter: natural-deduction
% Section: identity

\documentclass[../../../include/open-logic-section]{subfiles}

\begin{document}

\olfileid{fol}{ntd}{ide}

\olsection{\usetoken{S}{identity} கொண்ட \usetoken{P}{derivation}}

!!{identity} கொண்ட !!^{derivation}s க்குக் கூடுதல் அனுமான விதிகள் தேவை.

\begin{defish}
\AxiomC{}
\RightLabel{\Intro{\eq}}
\UnaryInfC{$\eq[t][t]$}
\DisplayProof
\hfill
\begin{tabular}{r}
\AxiomC{$\eq[t_1][t_2]$}
\AxiomC{$!A(t_1)$}
\RightLabel{\Elim{\eq}}
\BinaryInfC{$!A(t_2)$}
\DisplayProof
\\[3ex]
\AxiomC{$\eq[t_1][t_2]$}
\AxiomC{$!A(t_2)$}
\RightLabel{\Elim{\eq}}
\BinaryInfC{$!A(t_1)$}
\DisplayProof
\end{tabular}
\end{defish}

% SEGMENT OLP-0096-S02
மேலுள்ள விதிகளில் $t$, $t_1$, $t_2$ ஆகியவை மூடிய சொற்கள். எடுகோள்கள்
எவையும் இல்லாமல், $\eq[t][t]$ வடிவிலான எந்த அடையாளக் கூற்றையும் நேரடியாக
!!{derive} செய்ய \Intro{\eq} விதி அனுமதிக்கிறது.

\begin{ex}
$s$, $t$ ஆகியவை மூடிய சொற்கள் எனில், $!A(s), \eq[s][t] \Proves !A(t)$:
\begin{prooftree}
\AxiomC{$\eq[s][t]$}
\AxiomC{$!A(s)$}
\RightLabel{$\Elim{\eq}$}
\BinaryInfC{$!A(t)$}
\end{prooftree}
இது ``ஒத்தனவற்றின் பதிலிடத்தக்கதன்மைக் கொள்கை'' அல்லது லைப்னிட்சின் விதி
என்ற பெயர்களில் அறிமுகமானதாக இருக்கலாம்.
\end{ex}

% SEGMENT OLP-0096-S03
\begin{prob}
$=$ ஆனது சமச்சீரும் கடப்புடையதும் என்பதை நிறுவுக; அதாவது,
$\lforall[x][\lforall[y][(\eq[x][y] \lif
    \eq[y][x])]]$ மற்றும் $\lforall[x][\lforall[y][\lforall[z]((\eq[x][y]
    \land \eq[y][z]) \lif \eq[x][z])]]$ ஆகியவற்றின் !!{derivation}s ஐக்
கொடுக்கவும்.
\end{prob}

% SEGMENT OLP-0096-S04
\begin{ex}
பின்வரும் !!{sentence}[acc] நாம் !!{derive} செய்கிறோம்:
\begin{align*}
& \lforall[x][\lforall[y][((!A(x) \land !A(y)) \lif \eq[x][y])]]
\intertext{பின்வரும் !!{sentence} இலிருந்து:}
& \lexists[x][\lforall[y][(!A(y) \lif \eq[y][x])]]
\end{align*}
!!{derivation}[acc] பின்னோக்கி உருவாக்குகிறோம்:
\begin{prooftree}
\AxiomC{$\lexists[x][\lforall[y][(!A(y) \lif \eq[y][x])]]
\quad \Discharge{!A(a) \land !A(b)}{1}$}
\DeduceC{$\eq[a][b]$}
\DischargeRule{\Intro{\lif}}{1}
\UnaryInfC{$((!A(a) \land !A(b)) \lif \eq[a][b])$}
\RightLabel{\Intro{\lforall}}
\UnaryInfC{$\lforall[y][((!A(a) \land !A(y)) \lif \eq[a][y])]$}
\RightLabel{\Intro{\lforall}}
\UnaryInfC{$\lforall[x][\lforall[y][((!A(x) \land !A(y)) \lif \eq[x][y])]]$}
\end{prooftree}

% SEGMENT OLP-0096-S05
இப்போது முதன்மை எடுகோளைப் பயன்படுத்த வேண்டும். அது ஓர் இருத்தலளவை
!!{formula} என்பதால், இடைநிலை முடிவு~$\eq[a][b]$ ஐ !!{derive} செய்ய
\Elim{\lexists} ஐப் பயன்படுத்துகிறோம்.
\begin{prooftree}
\AxiomC{$\lexists[x][\lforall[y][(!A(y) \lif \eq[y][x])]]$}
\AxiomC{$\Discharge{\lforall[y][(!A(y) \lif \eq[y][c]])}{2}$}
\noLine
\UnaryInfC{$\Discharge{!A(a) \land !A(b)}{1}$}
\DeduceC{$\eq[a][b]$}
\DischargeRule{\Elim{\lexists}}{2}
\BinaryInfC{$\eq[a][b]$}
\DischargeRule{\Intro{\lif}}{1}
\UnaryInfC{$((!A(a) \land !A(b)) \lif \eq[a][b])$}
\RightLabel{\Intro{\lforall}}
\UnaryInfC{$\lforall[y][((!A(a) \land !A(y)) \lif \eq[a][y])]$}
\RightLabel{\Intro{\lforall}}
\UnaryInfC{$\lforall[x][\lforall[y][((!A(x) \land !A(y)) \lif \eq[x][y])]]$}
\end{prooftree}

% SEGMENT OLP-0096-S06
மேல் வலப்புறத் துணை-!!{derivation} ஐ நிறைவுசெய்ய, அதன் எடுகோள்களைப்
பயன்படுத்தி $\eq[a][c]$, $\eq[b][c]$ ஆகியவற்றைக் காட்டுகிறோம். இதற்கு இரு
தனித்த !!{derivation}s தேவை. $\eq[a][c]$ க்கான !!{derivation} பின்வருமாறு:
\begin{prooftree}
\AxiomC{$\Discharge{\lforall[y][(!A(y) \lif \eq[y][c]])}{2}$}
\RightLabel{\Elim{\lforall}}
\UnaryInfC{$!A(a) \lif \eq[a][c]$}
\AxiomC{$\Discharge{!A(a) \land !A(b)}{1}$}
\RightLabel{\Elim{\land}}
\UnaryInfC{$!A(a)$}
\RightLabel{\Elim{\lif}}
\BinaryInfC{$\eq[a][c]$}
\end{prooftree}
$\eq[a][c]$, $\eq[b][c]$ ஆகியவற்றிலிருந்து \Elim{\eq} ஆல் $\eq[a][b]$ ஐ
!!{derive} செய்கிறோம்.
\end{ex}

% SEGMENT OLP-0096-S07
\begin{prob}
பின்வரும் !!{formula}s[gen] !!{derivation}s ஐக் கொடுக்கவும்:
\begin{enumerate}
\item $\lforall[x][\lforall[y][((\eq[x][y] \land !A(x)) \lif !A(y))]]$
\item $\lexists[x][!A(x)] \land \lforall[y][\lforall[z][((!A(y) \land
    !A(z)) \lif \eq[y][z])]] \lif \lexists[x][(!A(x) \land
  \lforall[y][(!A(y) \lif \eq[y][x])])]$
\end{enumerate}
\end{prob}

\end{document}
